\subsection{Cross-ortholog selection-imprint audit for $B^{\ast}_{Q6}$ residuals}
\label{sec:bioreality-q6-cross-ortholog-selection-imprint}
\origin{ai}

\paragraph{Finite row surface.}
The selection-imprint packet is evaluated on a finite cross-ortholog table rather than on a free biological interpretation of the $B^{\ast}_{Q6}$ coordinates.  The registered run used OrthoDB v12 Fungi level 4751 single-copy orthogroups, Ensembl Fungi release-63 coding sequences, and PaxDb v6 integrated protein-abundance ppm as the abundance readout.  After the reconciliation filters, the tested surface contained $80$ orthogroups, $3$ species, and $240$ joined rows.  Each row carries a coding-sequence residual profile, an orthogroup label, a species label, and an abundance proxy.  The admissible question is therefore narrow: whether the $B^{\ast}_{Q6}$ residual coordinates add a stable cross-ortholog imprint beyond the ordinary codon-bias and expression-abundance structure visible in those rows.

\paragraph{Coordinate family and bounded translation context.}
The residual family available to the packet is the nine-coordinate $B^{\ast}_{Q6}$ basis
$\{\mathrm{Arg}_{AGR}, \mathrm{Ile}_{AUA}, K_{AAA}, \mathrm{Leu}_{CUN/UUR}, \mathrm{Leu}_{UUA/UUG}, \mathrm{Ser}_{UCA/UCG}, \mathrm{Ser}_{UCR/AGY}, \mathrm{Thr}_{ACR/ACY}, f_3\text{-stress}\}$.
A separate modeled-translation contact confirmed that these coordinates can be projected against tAI-style readouts: across three organisms, the modeled tAI matrix had $81$ entries above the bounded-contact threshold $0.01$, with maximum absolute coupling $3{,}801{,}961.031607896$.  That contact is useful only as a bounded codon-adaptation readback.  It is not an observation of ribosome motion, mature protein amount, folding, physical feasibility, biological function, or a universal rule.  The present ortholog audit therefore treats translation context as background discipline, not as automatic evidence for selection.

\paragraph{Selection-imprint statistic.}
The positive gate for a strong evolutionary selection-imprint claim was evaluated by a finite comparison against an effective null surface.  The measured improvement was
$$
\begin{aligned}
\Delta R^2 &= 0.03250515015613468,\\
\Delta L_{\mathrm{bits}} &= -258.6010084301837,\\
\Delta R^2_{\mathrm{leave\mbox{-}clade\mbox{-}out}} &= -0.0021477729853302424,\\
\max \Delta R^2_{\mathrm{null},95} &= 0.05427157066522958.
\end{aligned}
$$
The in-sample residual improvement is positive, but it is smaller than the recorded $95\%$ null envelope, and the leave-clade-out comparison is slightly negative.  Consequently the finite gate returns
$$
\begin{aligned}
\mathrm{positive\_selection\_imprint} &= \mathrm{false},\\
\mathrm{claim\_supported} &= \mathrm{false},\\
\mathrm{verdict} &= \text{codon-bias--expression correlation}.
\end{aligned}
$$
The status value ``passed'' means that the offline computation completed and produced this bounded verdict; it does not mean that the strong biological claim passed.

\paragraph{Relation to protein-abundance mediation.}
The negative imprint verdict is consistent with the broader protein-abundance discipline for the same residual family.  In a separate multi-organism mediation audit, per-protein modeled tAI and PaxDb abundance rows were joined for $12$ organisms, with joined protein counts ranging from $1{,}656$ in \emph{Sulfolobus solfataricus} to $19{,}395$ in mouse, and all organism-specific residual designs reported powered rank sufficiency.  Those rows support a statistical decomposition of $Q \to T \to P$ surfaces under controls, not a causal selection mechanism.  For example, the yeast mediation fraction was $0.9491551041216973$, while the human fraction was $0.27206289160119157$ and the mouse fraction was $0.15488799907323988$.  Such finite decompositions can show that a modeled translation covariate absorbs part of a protein-abundance association, but they do not turn the three-species ortholog table into a proof of evolutionary selection.

\paragraph{Axis specificity.}
A separate structural-order audit found exactly one significant coordinate under its own fold-level sign-flip gate: $f_3$-stress had weighted held-out $\Delta R^2 = 0.0031744502294756994$ and sign-flip right-tail probability $0.03125$.  The remaining coordinates did not pass that gate; for instance $\mathrm{Arg}_{AGR}$ had $p=0.142578125$, $\mathrm{Ile}_{AUA}$ had $p=0.0693359375$, and $\mathrm{Thr}_{ACR/ACY}$ had negative weighted held-out increment $-0.00045340330387383447$.  This axis-specific fact should not be imported as selection-imprint evidence.  It only says that one coordinate carried a structural-order signal in a different finite design.

\paragraph{Local cannot-claim boundary.}
The allowed conclusion is an honest bounded one: in this $80$-orthogroup, $240$-row, $3$-species table, the observed $B^{\ast}_{Q6}$ residual statistic is best reported as codon-bias--expression correlation rather than as a powered cross-ortholog selection imprint.  PaxDb ppm values are integrated across studies and species; they are not matched-condition absolute per-cell protein measurements.  Cross-reference reconciliation can remove proteins non-randomly, and a positive residual statistic would at most be an imprint criterion, not a direct molecular selection mechanism.  No translation realization, ribosome mechanism, folding claim, physical admissibility claim, molecular-function claim, organismal phenotype claim, or universal biological law follows from this table unless a separate layer-matched contact supplies that measurement and the corresponding controlled statistic survives its own null surface.